\documentclass[10pt,a4paper]{amsart}
\usepackage[margin=19mm]{geometry}
\usepackage{amsmath,amssymb,mathtools}
\usepackage[scr=rsfs,cal=euler]{mathalfa}
\usepackage{xfrac}
\usepackage[unicode,hidelinks,bookmarks=false]{hyperref}
\newcommand{\ie}{i.e.\ }
\newcommand{\cf}{cf.\ }
\newcommand{\resp}{respectively\ }
\newcommand{\Subsection}{\subsection}
\providecommand{\nobreakdash}{\nobreak\hbox{-}}
\setlength{\emergencystretch}{2em}
\multlinegap0pt
\usepackage{amssymb}
\newtheorem{proposition}{Proposition}
\newcommand{\R}{\mathbb{R}}
\newcommand{\Sph}{\mathbb{S}^{2}}
\DeclareMathOperator{\dist}{dist}
\title{The mean distance to a simple closed curve on the sphere}
\author{J. V. M. Pimentel}
\date{Source: arXiv:2609.09638v1 (CC BY 4.0)}
\begin{document}
\maketitle

\noindent\emph{Source and licence.} The mean distance to a simple closed curve on the sphere. Version arXiv:2609.09638v1 (CC BY 4.0), distributed by arXiv under the Creative Commons Attribution 4.0 International licence, http://creativecommons.org/licenses/by/4.0/, as recorded in the arXiv licence metadata for this exact version and shown on the abstract page https://arxiv.org/abs/2609.09638v1. Full licence notice: \texttt{licenses/arxiv-2609.09638-CC-BY-4.0.txt}.\par\smallskip

\noindent\textit{Editorial scope.} Counted unit: the article's proposition prop:short (source0.tex lines 97--108). The article's complete original proof of it is delivered with it (source0.tex lines 190--260, one \texttt{proof} environment of 71 lines, which the article places away from the statement). No mathematics is written, repaired or re-derived: the statement and the proof are the article's own lines, in the article's own notation, and no retained line is substituted or re-flowed.  No further source-local context is needed: every cross-reference of every retained line resolves inside the two delivered spans.  Nothing was omitted from any retained span, so the package carries no omission record.  No retained line contains a citation, so the wrapper carries no bibliography and no bibliography entry.  No retained line contains a figure environment, an image-inclusion command or drawing code, so no image is delivered and the package declares no asset dependency.  Editorial formatting only, in addition: an AMS \texttt{amsart} preamble that restates the article's own declarations and macros used by the retained lines, namely the environments \texttt{proposition} on separate counters, together with the standard packages those lines need.  The retained environments print in this document as Proposition 1 in order of appearance, whereas the article prints the same items with the numbering of its own full document.  The complete native source of the article as distributed in the arXiv e-print for this exact version is bundled unchanged as \texttt{sources/209840/source0.tex}.
\begin{proposition}\label{prop:short}
Every rectifiable closed curve $\mathcal{C}\subset\Sph$ of length $L$, simple or
not, satisfies
\begin{equation}\label{eq:short}
  \mathcal{J}(\mathcal{C})\ \ge\ \frac{\pi}{2}-\frac{L}{2\pi} .
\end{equation}
For $L\le2\pi$ the circles of length $L$ attain it, and no other simple closed
curve does, so
\begin{equation}
  J(L)=\frac{\pi}{2}-\frac{L}{2\pi},\qquad 0<L\le2\pi .
\end{equation}
\end{proposition}

\begin{proof}[Proof of Proposition~\ref{prop:short}]
Parametrize $\mathcal{C}$ by arclength as $\gamma\colon\R/L\mathbb{Z}\to\Sph$, fix
$x\in\Sph$ and put $\psi_{x}(s)=\dist(x,\gamma(s))$, a $1$-Lipschitz function.
Its minimum is $\dist(x,\mathcal{C})$ and its maximum is
$\pi-\dist(-x,\mathcal{C})$, because $\dist(-x,\,\cdot\,)=\pi-\dist(x,\,\cdot\,)$
turns the maximum of the one into the minimum of the other. A continuous
function on a circle has total variation at least twice its oscillation, so
\begin{equation}\label{eq:tv}
  \int_{0}^{L}\lvert\psi_{x}'(s)\rvert\,ds\ \ge\
  2\bigl(\pi-\dist(x,\mathcal{C})-\dist(-x,\mathcal{C})\bigr).
\end{equation}
At every $s$ where $\gamma'(s)$ exists and $x\ne\pm\gamma(s)$ the chain rule
applies to $\psi_{x}(s)=\arccos\langle x,\gamma(s)\rangle$ and gives
$\lvert\psi_{x}'(s)\rvert=\lvert\cos\varphi\rvert$, where $\varphi$ is the azimuth
of $x$ in geodesic polar coordinates centered at $\gamma(s)$, measured from
$\gamma'(s)$. Both exceptional sets are null, the second because $\gamma'$ would
vanish at a density point of $\gamma^{-1}(\{\pm x\})$. Integrating,
\[
  \int_{\Sph}\lvert\cos\varphi\rvert\,dA(x)
  =\int_{0}^{\pi}\!\!\int_{0}^{2\pi}\lvert\cos\varphi\rvert\,d\varphi\,\sin\theta\,d\theta
  =8 ,
\]
so Fubini gives $\int_{\Sph}\int_{0}^{L}\lvert\psi_{x}'(s)\rvert\,ds\,dA(x)=8L$. The
antipodal map preserves $dA$, so each of the two distance integrals equals
$4\pi\mathcal{J}(\mathcal{C})$, and integrating \eqref{eq:tv} over $\Sph$ gives
$8L\ge8\pi^{2}-16\pi\mathcal{J}(\mathcal{C})$, which is \eqref{eq:short}.

Let now $L\le2\pi$ and let $\mathcal{C}$ be the circle of geodesic radius
$\theta_{0}=\arcsin(L/2\pi)\le\pi/2$ about a fixed center. A point at geodesic
distance $\theta$ from the center is at distance $\lvert\theta-\theta_{0}\rvert$
from $\mathcal{C}$, so
\[
  4\pi \mathcal{J}(\mathcal{C})
  =2\pi\int_{0}^{\pi}\lvert\theta-\theta_{0}\rvert\sin\theta\,d\theta
  =2\pi(\pi-2\sin\theta_{0}),
\]
that is $\mathcal{J}(\mathcal{C})=\pi/2-\sin\theta_{0}=\pi/2-L/(2\pi)$.

Finally, let $\mathcal{C}$ be a curve of length $L\le2\pi$ with
$\mathcal{J}(\mathcal{C})=\pi/2-L/(2\pi)$. Every step above except \eqref{eq:tv}
was an identity, so \eqref{eq:tv} is an equality for almost every $x$. For such
an $x$, each of the two arcs cut from the parameter circle by a minimum and a
maximum point of $\psi_{x}$ carries variation exactly equal to the oscillation,
so $\psi_{x}$ is monotone on both. A level strictly between the extremes is
therefore attained on at most one subinterval of each arc, and only countably
many levels are attained on a nondegenerate one. So, writing
$B(x,r)=\{\dist(\,\cdot\,,x)<r\}$ for the open geodesic ball, for almost every
$x$ and almost every $r$ the circle $\partial B(x,r)$ meets $\mathcal{C}$ in at
most two points.

Suppose $\mathcal{C}$ were not contained in a circle. Any three of its points
determine a circle, which some fourth point of $\mathcal{C}$ misses, so
$\mathcal{C}$ has four points that are not concyclic, hence not coplanar.
Label them $y_{1},y_{2},y_{3},y_{4}$ in the cyclic order in which they occur
along $\mathcal{C}$. Two line segments of $\R^{3}$ that meet have coplanar
endpoints, so $[y_{1},y_{3}]$ and $[y_{2},y_{4}]$ are disjoint compact convex
sets and some plane separates them strictly. As the $y_{i}$ lie on $\Sph$, that
plane can be written $\{\langle\,\cdot\,,x\rangle=\cos r\}$ with $x\in\Sph$,
$r\in(0,\pi)$ and $y_{1},y_{3}$ on the side of $x$; then
$y_{1},y_{3}\in B(x,r)$ while $y_{2},y_{4}\notin\overline{B(x,r)}$. Each of the
four arcs of $\mathcal{C}$ between consecutive $y_{i}$ joins a point of
$B(x,r)$ to a point outside $\overline{B(x,r)}$, hence meets
$\partial B(x,r)$; as $\mathcal{C}$ is simple these four arcs meet only in the
$y_{i}$, which lie off $\partial B(x,r)$, so $\partial B(x,r)$ meets
$\mathcal{C}$ in at least four points. The four inclusions are strict, so the
same holds on a neighborhood of $(x,r)$, contradicting the previous paragraph.

So $\mathcal{C}$ lies on a circle, and equals it because a proper closed subset
of a circle is not homeomorphic to one. Its length is $L$, so it is a circle of
geodesic radius $\theta_{0}$.
\end{proof}

\end{document}
